\documentclass[conference]{IEEEtran}

\usepackage{graphicx}
\usepackage{hyperref}

\title{News Credibility Analysis Using Logistic Regression}

\author{
\IEEEauthorblockN{Team X}
\IEEEauthorblockA{
Advik Khandelwal, Sahil Kumar, Harsh Ahalawat, Meet Kumar
}
}

\begin{document}

\maketitle

\begin{abstract}
This paper presents a machine learning based system for detecting
fake and real news articles using Natural Language Processing (NLP).
Text preprocessing, TF-IDF feature extraction, and Logistic Regression
classification are applied to analyze news content. The proposed model
achieves high accuracy and demonstrates an effective approach for
automatic misinformation detection.
\end{abstract}

\begin{IEEEkeywords}
Fake News Detection, Machine Learning, NLP, TF-IDF, Logistic Regression
\end{IEEEkeywords}

\section{Introduction}

The rapid growth of online media has increased the spread of
misinformation. Manual verification of news is difficult and
time-consuming. Machine learning techniques enable automated
credibility analysis by learning patterns from textual data.
This project develops a system that classifies news articles
as real or fake using supervised learning.

\section{Related Work}

Previous research in fake news detection uses Natural Language
Processing and classification algorithms such as Naive Bayes,
Support Vector Machines, and deep learning models. Traditional
machine learning methods remain effective due to lower computational
cost and strong performance on textual datasets.

\section{Methodology}

\subsection{Dataset Description}

The dataset was obtained using KaggleHub and contains news
articles with title, text, and label fields. The title and
article text were combined to create the input content.
Missing values were removed before training.

\subsection{Preprocessing}

Text preprocessing included:
\begin{itemize}
\item Lowercasing text
\item Removing URLs and special characters
\item Retaining alphanumeric words
\item Stopword removal using NLTK
\end{itemize}

\subsection{Feature Engineering}

TF-IDF vectorization was used with:
\begin{itemize}
\item Maximum features: 20,000
\item N-grams: (1,2)
\item Minimum document frequency: 5
\item Maximum document frequency: 0.7
\end{itemize}

\subsection{Model Used}

A Logistic Regression classifier was trained with:
\begin{itemize}
\item Max iterations: 2000
\item Class weight: balanced
\item Regularization parameter C = 2.0
\end{itemize}

\subsection{Training and Validation}

The dataset was split using an 80:20 train-test ratio with
stratified sampling. Model performance was evaluated using
precision, recall, F1-score, and confusion matrix metrics.

\section{Results}

\subsection{Quantitative Results}

\begin{verbatim}
Accuracy: 0.96

Classification Report:
Precision Recall F1
Class 0: 0.97 0.96 0.96
Class 1: 0.96 0.97 0.97
\end{verbatim}

Confusion Matrix:
\begin{verbatim}
[[6726  280]
 [ 232 7070]]
\end{verbatim}

\section{Discussion}

Initially, the model showed bias toward predicting news as fake
due to class imbalance in the dataset. This issue was addressed
using resampling techniques and balanced class weights during
training, which improved prediction fairness and performance.

\section{Conclusion}

The project demonstrates that Logistic Regression combined with
TF-IDF features provides an efficient and accurate solution for
fake news detection. The system achieves approximately 96\%
accuracy and supports real-time predictions through a Streamlit
application.

\section{Appendix}

Repository Link:

\begin{center}
\url{https://github.com/Advikkhandelwal/News_credibilty_analysis}
\end{center}

\end{document}